%% Acta de constitución del proyecto con cliente (plantilla)
%% Compilar: tectonic -X compile acta-constitucion-proyecto.tex
%%
%% No confundir con el acta de constitución de la consultora
%% (docs/institucional/anexo-acta-constitucion.tex, LGS-INST-002), que ya está
%% firmada y autoriza el inicio de actividades de la empresa. Esta plantilla
%% autoriza el inicio de un proyecto concreto, una vez elegido el cliente.
\documentclass[11pt,a4paper]{article}
\usepackage{../latex/legasoft}

\lgtitulo{Acta de constitución del proyecto}
\lgsubtitulo{Autorización formal de inicio y marco de referencia del proyecto}
\lgtituloCorto{Acta de constitución del proyecto}
\lgcodigo{LGS-PRY-001}
\lgversion{0.1}
\lgestado{Plantilla}
\lgfecha{\campo{fecha de aprobación}}
\lgautor{\lgDirector\ (\lgDirectorRol)}

\begin{document}
\lgportada

\lgcontrolversiones{
0.1 & \campo{dd-mm-aaaa} & \lgDirector & Plantilla inicial sin datos de proyecto. \\
}

\begin{porcompletar}[Cómo usar esta plantilla]
Esta plantilla fija la estructura del acta con que se autoriza el inicio de un
proyecto con cliente. Los campos marcados en ámbar deben sustituirse por los
datos reales una vez seleccionada la organización cliente; al completarla,
cambie el estado a \emph{Aprobado} y suba la versión a 1.0. El inicio de
actividades de la consultora ya está autorizado por el acta
\texttt{LGS-INST-002}, incluida como Anexo B del documento de entrega
\texttt{LGS-ENT-001}.
\end{porcompletar}

\tableofcontents
\clearpage

% =============================================================================
\section{Identificación del proyecto}

\begin{center}
\begin{tabularx}{\linewidth}{@{}L{4.6cm} Y@{}}
\toprule
\sffamily\bfseries\color{lgPrimario}Campo &
\sffamily\bfseries\color{lgPrimario}Contenido \\
\midrule
Nombre del proyecto        & \campo{nombre} \\
Código interno             & \campo{LGS-PRY-000} \\
Organización cliente       & \campo{razón social} \\
Patrocinador               & \campo{nombre y cargo} \\
Director de proyecto       & \lgDirector \\
Fecha de inicio prevista   & \campo{dd-mm-aaaa} \\
Fecha de cierre prevista   & \campo{dd-mm-aaaa} \\
Consultora ejecutora       & Legasoft \\
\bottomrule
\end{tabularx}
\captionof{table}{Datos de identificación del proyecto.}
\end{center}

% =============================================================================
\section{Justificación y necesidad}

Descripción del problema u oportunidad que motiva el proyecto, sustentada en el
diagnóstico realizado sobre el funcionamiento actual de la organización cliente.

\begin{porcompletar}
\campo{Situación actual observada: procesos manuales, registros dispersos,
ineficiencias, restricciones o sistemas que ya no responden a los procesos.}

\campo{Consecuencias medibles de esa situación: tiempos, errores, costos,
riesgos o pérdida de información.}
\end{porcompletar}

% =============================================================================
\section{Objetivos del proyecto}

\subsection{Objetivo general}

\begin{porcompletar}
\campo{Un enunciado que exprese el resultado esperado del proyecto en términos
del beneficio para el cliente, no de la tecnología empleada.}
\end{porcompletar}

\subsection{Objetivos específicos}

\begin{porcompletar}
\campo{Entre tres y seis objetivos verificables, cada uno asociado a un
entregable identificable.}
\end{porcompletar}

% =============================================================================
\section{Alcance}

\subsection{Incluido en el alcance}

\begin{porcompletar}
\campo{Procesos, módulos, integraciones y entregables comprendidos en el
proyecto.}
\end{porcompletar}

\subsection{Excluido del alcance}

Declarar explícitamente lo que el proyecto no cubre evita expectativas que no
puedan sostenerse, en coherencia con el valor de transparencia de Legasoft.

\begin{porcompletar}
\campo{Funcionalidades, sistemas o servicios que quedan fuera y, cuando
corresponda, la razón.}
\end{porcompletar}

% =============================================================================
\section{Entregables principales}

\begin{center}
\begin{tabularx}{\linewidth}{@{}C{2.0cm} Y L{4.0cm} C{2.4cm}@{}}
\toprule
\sffamily\bfseries\color{lgPrimario}Código &
\sffamily\bfseries\color{lgPrimario}Entregable &
\sffamily\bfseries\color{lgPrimario}Responsable &
\sffamily\bfseries\color{lgPrimario}Fecha \\
\midrule
ENT-01 & Especificación de requisitos de software (SRS) & \lgDirectorRol & \campo{fecha} \\
ENT-02 & Documento de arquitectura y modelos C4        & \lgArquitectoRol & \campo{fecha} \\
ENT-03 & Estrategia de calidad y plan de pruebas       & \lgQARol & \campo{fecha} \\
ENT-04 & \campo{entregable}                            & \campo{rol} & \campo{fecha} \\
\bottomrule
\end{tabularx}
\captionof{table}{Entregables comprometidos y sus responsables.}
\end{center}

% =============================================================================
\section{Hitos}

\begin{center}
\begin{tabularx}{\linewidth}{@{}C{2.0cm} Y C{2.6cm}@{}}
\toprule
\sffamily\bfseries\color{lgPrimario}Hito &
\sffamily\bfseries\color{lgPrimario}Descripción &
\sffamily\bfseries\color{lgPrimario}Fecha prevista \\
\midrule
H1 & Cierre del levantamiento de requisitos          & \campo{fecha} \\
H2 & Aprobación de la arquitectura                   & \campo{fecha} \\
H3 & Entrega funcional para pruebas de aceptación    & \campo{fecha} \\
H4 & Entrega final y cierre del proyecto             & \campo{fecha} \\
\bottomrule
\end{tabularx}
\captionof{table}{Hitos de control del proyecto.}
\end{center}

% =============================================================================
\section{Interesados}

\begin{center}
\begin{tabularx}{\linewidth}{@{}L{3.8cm} L{3.4cm} Y@{}}
\toprule
\sffamily\bfseries\color{lgPrimario}Interesado &
\sffamily\bfseries\color{lgPrimario}Rol &
\sffamily\bfseries\color{lgPrimario}Interés o expectativa \\
\midrule
\campo{nombre} & Patrocinador        & \campo{expectativa} \\
\campo{nombre} & Usuario clave       & \campo{expectativa} \\
\campo{nombre} & Responsable técnico & \campo{expectativa} \\
\bottomrule
\end{tabularx}
\captionof{table}{Registro inicial de interesados.}
\end{center}

% =============================================================================
\section{Equipo del proyecto y responsabilidades}

La asignación de roles corresponde a la estructura definida en el documento de
constitución de Legasoft (\texttt{LGS-INST-001}, sección 10).

\begin{center}
\begin{tabularx}{\linewidth}{@{}L{4.4cm} L{4.0cm} Y@{}}
\toprule
\sffamily\bfseries\color{lgPrimario}Integrante &
\sffamily\bfseries\color{lgPrimario}Rol &
\sffamily\bfseries\color{lgPrimario}Responsabilidad en este proyecto \\
\midrule
\lgDirector   & \lgDirectorRol   & Coordinación, requisitos, seguimiento y relación con el cliente. \\
\lgArquitecto & \lgArquitectoRol & Arquitectura, modelo de datos, APIs y lógica de negocio. \\
\lgQA         & \lgQARol         & Interfaz de usuario, estrategia de calidad y validación de entregas. \\
\bottomrule
\end{tabularx}
\captionof{table}{Equipo asignado al proyecto.}
\end{center}

% =============================================================================
\section{Supuestos y restricciones}

\subsection{Supuestos}

\begin{porcompletar}
\campo{Condiciones que se dan por ciertas para planificar: disponibilidad de
usuarios para entrevistas, acceso a datos, existencia de infraestructura, etc.}
\end{porcompletar}

\subsection{Restricciones}

\begin{porcompletar}
\campo{Límites de tiempo, presupuesto, tecnología obligatoria, normativa
aplicable o capacidad operativa del cliente.}
\end{porcompletar}

% =============================================================================
\section{Riesgos iniciales}

\begin{center}
\begin{tabularx}{\linewidth}{@{}C{1.6cm} Y C{1.8cm} C{1.8cm} L{3.6cm}@{}}
\toprule
\sffamily\bfseries\color{lgPrimario}ID &
\sffamily\bfseries\color{lgPrimario}Riesgo &
\sffamily\bfseries\color{lgPrimario}Prob. &
\sffamily\bfseries\color{lgPrimario}Impacto &
\sffamily\bfseries\color{lgPrimario}Respuesta prevista \\
\midrule
R-01 & \campo{descripción} & \campo{A/M/B} & \campo{A/M/B} & \campo{acción} \\
R-02 & \campo{descripción} & \campo{A/M/B} & \campo{A/M/B} & \campo{acción} \\
R-03 & \campo{descripción} & \campo{A/M/B} & \campo{A/M/B} & \campo{acción} \\
\bottomrule
\end{tabularx}
\captionof{table}{Riesgos identificados al inicio del proyecto.}
\end{center}

% =============================================================================
\section{Criterios de éxito y de aceptación}

\begin{porcompletar}
\campo{Condiciones observables que permitirán declarar el proyecto exitoso:
requisitos cubiertos, criterios de aceptación aprobados, atributos de calidad
verificados según ISO/IEC 25010 y compromisos de plazo cumplidos.}
\end{porcompletar}

% =============================================================================
\section{Aprobación}

Con la firma de este documento se autoriza formalmente el inicio del proyecto y
se reconoce la asignación de los recursos descritos.

\vspace{1.2cm}

\begin{center}
\begin{tabularx}{\linewidth}{@{}Y Y Y@{}}
\centering\rule{4.6cm}{0.5pt} &
\centering\rule{4.6cm}{0.5pt} &
\centering\arraybackslash\rule{4.6cm}{0.5pt} \\[-2pt]
\centering\sffamily\small\campo{patrocinador} &
\centering\sffamily\small\lgDirector &
\centering\arraybackslash\sffamily\small\campo{representante cliente} \\[-4pt]
\centering\scriptsize\color{lgGris}Patrocinador &
\centering\scriptsize\color{lgGris}Director de Proyecto &
\centering\arraybackslash\scriptsize\color{lgGris}Por la organización cliente \\
\end{tabularx}
\end{center}

\end{document}
